\chapter{Topologie van metrische ruimten}\label{ch:b2:metric}

De topologie van de reële rechte (volume van bachelorjaar 1)
veralgemeent zich, vrijwel zonder één woord te veranderen, tot elke
verzameling die van een afstand is voorzien. De winst is enorm:
rijen van functies, matrices, krommen --- alles wordt een punt van
een \omterm{def:b2:metric:def}{metrische ruimte}, en de drie pijlers die we hier bewijzen ---
\omterm{def:b2:metric:complete}{volledigheid} met de vastepuntsstelling van Banach, \omterm{def:b2:metric:compact}{compactheid},
\omterm{def:b2:metric:connected}{samenhang} --- gelden er eenvormig voor. Dit hoofdstuk is de
ruggengraat van de hele analysehelft van het boek.

\section{Metrische ruimten}

\begin{definition}\label{def:b2:metric:def}
Een \emph{metrische ruimte}\index{metrische ruimte} is een
verzameling $X$ met een afbeelding $d \colon X \times X \to \R_+$
zodanig dat voor alle $x, y, z$
\[
d(x,y) = 0 \iff x = y,
\qquad
d(x,y) = d(y,x),
\qquad
d(x,z) \leq d(x,y) + d(y,z).
\]
Bollen: $B(a, r) = \{x : d(a,x) < r\}$ (\omterm{def:b2:metric:topology}{open}), $\overline B(a,r) =
\{x : d(a,x) \leq r\}$ (gesloten). Een deelverzameling $A \subseteq
X$ wordt met de geïnduceerde afstand zelf een metrische ruimte.
\end{definition}

\begin{example}\label{ex:b2:metric:examples}
$\R$ met $\abs{x - y}$; $\R^n$ met elk van
\[
d_1(x,y) = \sum_i \abs{x_i - y_i},
\quad
d_2(x,y) = \Bigl(\sum_i (x_i - y_i)^2\Bigr)^{1/2},
\quad
d_\infty(x,y) = \max_i \abs{x_i - y_i};
\]
de verzameling $C(\intcc{a}{b})$ van de \omterm{def:b2:metric:continuity}{continue} functies met de
\emph{supafstand} $d_\infty(f, g) = \sup_{\intcc{a}{b}} \abs{f -
g}$ (eindig, want $f - g$ is begrensd); en elke verzameling met de
\emph{discrete} afstand ($d(x,y) = 1$ voor $x \neq y$). Afstanden
die uit normen voortkomen zijn het onderwerp van
\cref{ch:b2:nvs}.
\end{example}

\begin{definition}[Topologie van een metrische ruimte]\label{def:b2:metric:topology}
$U \subseteq X$ heet \emph{open}\index{open verzameling!in een
metrische ruimte} wanneer elk punt van $U$ het middelpunt is van
een bol die in $U$ ligt; $F$ heet \emph{gesloten} wanneer haar
complement open is. Omgevingen, inwendige, afsluiting, dichtheid en
rand worden precies gedefinieerd als op de reële rechte (volume van
bachelorjaar 1), met bollen in plaats van intervallen, en de daar
bewezen uitspraken --- verenigingen en doorsneden van open
verzamelingen, karakteriseringen van inwendige en afsluiting, de
afsluiting als kleinste gesloten bovenverzameling --- gaan met
dezelfde bewijzen over. Open bollen zijn open en gesloten bollen
zijn gesloten (driehoeksongelijkheid).
\end{definition}

\begin{example}[Inwendige, afsluiting en rand op één verzameling]\label{ex:b2:metric:intclosure}
Zij in $\R$ de verzameling $A = \intoc{0}{1} \cup \{2\}$.
Inwendige: $\intoo{0}{1}$ --- rond elke $x \in \intoo01$ blijft een
kleine bol in $A$; rond $1$ lekt elke bol $\intoo{1-r}{1+r}$ aan de
rechterkant uit $A$ weg, dus $1$ is geen inwendig punt; en het
geïsoleerde punt $2$ evenmin. Afsluiting: $\intcc{0}{1} \cup \{2\}$
(het punt $0$ is een limiet van $A$, verder komt er niets bij).
Rand (afsluiting min inwendige): $\{0, 1, 2\}$. Let op de
asymmetrieën die het onthouden waard zijn: een randpunt kan tot een
verzameling behoren zonder inwendig te zijn ($1$), het kan
aanhechtingspunt zijn zonder erbij te horen ($0$), en een
geïsoleerd punt is zijn eigen rand ($2$). Dezelfde boekhouding
loopt woordelijk in elke \omterm{def:b2:metric:def}{metrische ruimte}, met bollen in plaats van
intervallen.
\end{example}

\begin{definition}[Limieten, continuïteit]\label{def:b2:metric:continuity}
$x_n \to x$ in $X$ wanneer $d(x_n, x) \to 0$. Een afbeelding $f
\colon X \to Y$ tussen \omterm{def:b2:metric:def}{metrische ruimten} heet
\emph{continu}\index{continu} in $a$ wanneer
\[
\forall \varepsilon > 0,\ \exists\delta > 0,\quad
d_X(x, a) \leq \delta \implies d_Y\bigl(f(x), f(a)\bigr) \leq
\varepsilon ;
\]
gelijkwaardig (met hetzelfde bewijs als op $\R$): $f(x_n) \to f(a)$
voor elke rij $x_n \to a$. En $f$ heet
\emph{Lipschitz}\index{Lipschitz} met constante $k$ wanneer altijd
$d_Y(f(x), f(y)) \leq k\, d_X(x, y)$ --- dan is $f$ uniform
continu, en dus continu.
\end{definition}

\begin{theorem}[Globale karakterisering van continuïteit]\label{thm:b2:metric:globalcontinuity}
$f \colon X \to Y$ is \omterm{def:b2:metric:continuity}{continu} (in elk punt) dan en slechts dan als
het origineel van elke \omterm{def:b2:metric:topology}{open} verzameling \omterm{def:b2:metric:topology}{open} is --- dan en slechts
dan als het origineel van elke gesloten verzameling gesloten is.
\end{theorem}

\begin{proof}
($\Rightarrow$) Zij $V \subseteq Y$ \omterm{def:b2:metric:topology}{open} en $a \in f^{-1}(V)$: dan
ligt een bol $B(f(a), \varepsilon) \subseteq V$; de \omterm{def:b2:metric:continuity}{continuïteit} in
$a$ levert een $\delta$ met $f(B(a, \delta)) \subseteq B(f(a),
\varepsilon)$, dus $B(a, \delta) \subseteq f^{-1}(V)$.

($\Leftarrow$) Gegeven $a$ en $\varepsilon$: de verzameling
$f^{-1}\bigl(B(f(a), \varepsilon)\bigr)$ is \omterm{def:b2:metric:topology}{open} en bevat $a$, dus
bevat zij een bol $B(a, \delta)$: dat is juist de definitie van
\omterm{def:b2:metric:continuity}{continuïteit} in $a$. Voor gesloten verzamelingen: neem complementen
(\cref{prop:b2:structures:images}).
\end{proof}

\section{Volledige ruimten}

\begin{definition}\label{def:b2:metric:complete}
Een rij $(x_n)$ heet \emph{cauchyrij} wanneer $\sup_{p, q \geq N}
d(x_p, x_q) \to 0$ als $N \to \infty$. Een \omterm{def:b2:metric:def}{metrische ruimte} heet
\emph{volledig}\index{volledige metrische ruimte} wanneer elke
cauchyrij convergeert. Convergent $\Rightarrow$ cauchy geldt
altijd; gesloten deelverzamelingen van volledige ruimten zijn
volledig, en volledige deelverzamelingen van een willekeurige
ruimte zijn gesloten (met dezelfde bewijzen als op $\R$: volume van
bachelorjaar 1).
\end{definition}

\begin{example}[Cauchy zonder limiet]\label{ex:b2:metric:cauchynolimit}
Zij $X = \Q$ met de gebruikelijke afstand; de decimale
afkappingen van $\sqrt2$,
\[
x_0 = 1,\quad x_1 = 1.4,\quad x_2 = 1.41,\quad x_3 = 1.414,
\quad\dots
\]
voldoen aan $\abs{x_p - x_q} \leq 10^{-\min(p,q)}$: het is een
cauchyrij in $\Q$. Een limiet in $\Q$ zou ook de limiet in $\R$
zijn, namelijk $\sqrt2 \notin \Q$: in $X$ bestaat dus geen limiet.
Onvolledigheid is de aanwezigheid van zulke ``spooklimieten''; de
\omterm{def:b2:metric:complete}{volledigheid} van $\R$ werd in het volume van bachelorjaar 1 juist
zo geconstrueerd dat elke cauchyrij een thuis krijgt.
\end{example}

\begin{theorem}\label{thm:b2:metric:rncomplete}
$\R^n$ (met elk van de drie afstanden uit
\cref{ex:b2:metric:examples}) en $\bigl(C(\intcc{a}{b}),
d_\infty\bigr)$ zijn \omterm{def:b2:metric:complete}{volledig}.
\end{theorem}

\begin{proof}
$\R^n$: een cauchyrij is cauchy in elke coördinaat (voor alle drie
de afstanden geldt $\abs{x_i - y_i} \leq d(x,y)$), dus convergeert
elke coördinaat (\omterm{def:b2:metric:complete}{volledigheid} van $\R$, volume van bachelorjaar 1),
en coördinaatsgewijze convergentie impliceert convergentie voor
$d_\infty$ (eindig veel coördinaten), en dus voor alle drie (de
drie afstanden domineren elkaar op constante factoren na:
$d_\infty \leq d_2 \leq d_1 \leq n\,d_\infty$).

$C(\intcc{a}{b})$: zij $(f_n)$ een cauchyrij voor $d_\infty$. Voor
elke $x$ is $(f_n(x))$ een cauchyrij in $\R$ ($\abs{f_p(x) -
f_q(x)} \leq d_\infty(f_p, f_q)$) en convergeert dus naar zekere
$f(x)$. Laat in $\abs{f_p(x) - f_q(x)} \leq \varepsilon$ (geldig
voor $p, q \geq N_\varepsilon$ en alle $x$) de index $q \to \infty$
gaan: $\abs{f_p(x) - f(x)} \leq \varepsilon$ voor alle $x$, dat wil
zeggen $d_\infty(f_p, f) \leq \varepsilon$: uniforme convergentie.
De limiet is \omterm{def:b2:metric:continuity}{continu}: kies bij gegeven $\varepsilon$ een $p$ met
$\sup\abs{f_p - f} \leq \varepsilon$, en gebruik dan de \omterm{def:b2:metric:continuity}{continuïteit}
van $f_p$ in $a$ samen met de splitsing in drie termen
\[
\abs{f(x) - f(a)} \leq \abs{f(x) - f_p(x)} + \abs{f_p(x) - f_p(a)}
+ \abs{f_p(a) - f(a)} \leq 3\varepsilon
\]
voor $x$ dicht bij $a$. (Dit ``$3\varepsilon$-argument'' keert
terug als de stelling over uniforme limieten in
\cref{ch:b2:funcseq}.)
\end{proof}

\begin{example}[Open en gesloten verzamelingen herkend via continuïteit]\label{ex:b2:metric:recognize}
De globale karakterisering
(\cref{thm:b2:metric:globalcontinuity}) is het dagelijkse
gereedschap voor topologische boekhouding. In $\R^2$ is de
verzameling $\{(x, y) : x^2 + y^2 < 1,\ y > x^3\}$ \omterm{def:b2:metric:topology}{open} --- zij is
$g^{-1}(\intoo{-\infty}{1}) \cap h^{-1}(\intoo{0}{+\infty})$ voor
de \omterm{def:b2:metric:continuity}{continue} $g(x,y) = x^2 + y^2$ en $h(x, y) = y - x^3$, een
doorsnede van twee \omterm{def:b2:metric:topology}{open} originelen. In $\bigl(C(\intcc01),
d_\infty\bigr)$ is de verzameling van de functies met $f(0) = f(1)$
en $\int_0^1 f = 0$ gesloten --- het origineel van $\{(0,0)\}$
onder de \omterm{def:b2:metric:continuity}{continue} afbeelding $f \mapsto \bigl(f(0) - f(1),\
\int_0^1 f\bigr)$ naar $\R^2$ (elke coördinaat is
$1$-Lipschitz, zoals in \cref{exo:b2:metric:3}). De methode tekent
nooit een plaatje: geef een \omterm{def:b2:metric:continuity}{continue} afbeelding, lees de
verzameling als origineel, en beroep je op de stelling.
\end{example}

\begin{example}[Een gesloten verzameling met oneindig veel voorwaarden]\label{ex:b2:metric:lipclosed}
In $\bigl(C(\intcc01), d_\infty\bigr)$ is de verzameling
\[
L = \{f : \abs{f(x) - f(y)} \leq \abs{x - y}
\ \text{voor alle } x, y\}
\]
van de $1$-Lipschitz-functies gesloten, hoewel zij door
overaftelbaar veel voorwaarden wordt uitgesneden: voor elk vast
paar $(x, y)$ is de afbeelding $f \mapsto \abs{f(x) - f(y)} -
\abs{x - y}$ \omterm{def:b2:metric:continuity}{continu} (evaluaties zijn $1$-Lipschitz), dus
definieert elke afzonderlijke voorwaarde een gesloten verzameling,
en $L$ is de \emph{doorsnede} van deze familie --- en een
willekeurige doorsnede van gesloten verzamelingen is gesloten.
Hetzelfde sjabloon bewijst de geslotenheid voor monotone functies,
convexe functies of functies die door een vaste $g$ worden begrensd:
uniforme limieten erven elke eigenschap die als een familie van
gesloten puntsgewijze voorwaarden te schrijven is. Wat uniforme
limieten \emph{niet} vanzelf erven --- differentieerbaarheid
bijvoorbeeld --- is juist waarvoor \cref{ch:b2:funcseq} moet
zwoegen.
\end{example}

\begin{theorem}[Vastepuntsstelling van Banach]\label{thm:b2:metric:banach}
Zij $X$ een niet-lege \omterm{def:b2:metric:complete}{volledige metrische ruimte} en $f \colon X \to
X$ een \emph{contractie}: \omterm{def:b2:metric:continuity}{Lipschitz} met constante $k < 1$. Dan
heeft $f$ precies één vast punt $\ell$, en elke baan $x_{n+1} =
f(x_n)$ convergeert naar $\ell$, met\index{vastepuntsstelling van
Banach}
\[
d(x_n, \ell) \leq \frac{k^n}{1 - k}\, d(x_1, x_0) .
\]
\end{theorem}

\begin{proof}
Eenduidigheid: twee vaste punten liggen op een afstand die $\leq k$
maal zichzelf is. Bestaan: met inductie is $d(x_{n+1}, x_n) \leq
k^n d(x_1, x_0)$, dus voor $q > p$
\[
d(x_q, x_p) \leq \sum_{j=p}^{q-1} d(x_{j+1}, x_j)
\leq d(x_1, x_0) \sum_{j \geq p} k^j
= \frac{k^p}{1-k}\, d(x_1, x_0) \xrightarrow[p\to\infty]{} 0 :
\]
een cauchyrij, die dus convergeert naar zekere $\ell$; de
\omterm{def:b2:metric:continuity}{continuïteit} van $f$ laat $x_{n+1} = f(x_n)$ naar de limiet
overgaan: $\ell = f(\ell)$. De foutgrens is dezelfde schatting met
$q \to \infty$.
\end{proof}

\begin{example}[Een integraalvergelijking]\label{ex:b2:metric:integralequation}
Beschouw op $X = C(\intcc{0}{1})$ (\omterm{def:b2:metric:complete}{volledig},
\cref{thm:b2:metric:rncomplete}) de afbeelding $T(f)(x) = 1 +
\frac12 \int_0^x f(t)\,\dd t$. Voor $f, g \in X$ is
\[
\abs{T(f)(x) - T(g)(x)} \leq \frac12 \int_0^x \abs{f - g} \leq
\frac12\, d_\infty(f, g),
\]
dus is $T$ een $\frac12$-contractie: zij heeft precies één \omterm{def:b2:metric:continuity}{continu}
vast punt --- de oplossing van $f' = \frac f2$ met $f(0) = 1$,
namelijk $\eu^{x/2}$. Dit schema wordt, geïndustrialiseerd, de
stelling van Cauchy--Lipschitz in \cref{ch:b2:diffeq}.
\end{example}

\begin{example}[Een numeriek vast punt: $x = \cos x$]\label{ex:b2:metric:cosfix}
Op de volledige ruimte $X = \intcc{0}{1}$ beeldt $f = \cos$ de
ruimte $X$ in $\intcc{\cos 1}{1} \subseteq X$ af, en is zij een
contractie: volgens de middelwaardeongelijkheid geldt
\[
\abs{\cos x - \cos y} \leq \bigl(\sup_{\intcc01}\abs{\sin}\bigr)
\abs{x - y} = (\sin 1)\abs{x - y},
\qquad \sin 1 \approx 0.841 < 1 .
\]
Banach: er is precies één oplossing van $x = \cos x$ in
$\intcc01$ (en dus in $\R$: elk reëel vast punt ligt in
$\intcc{-1}{1}$ en, na één toepassing, in $\intcc{\cos 1}{1}$), en
de iteratie $x_{n+1} = \cos x_n$ convergeert er vanuit elk
beginpunt naar: $x_\infty \approx 0.739085$, het beroemde getal dat
men krijgt door op een rekenmachine op de cosinustoets te blijven
rammen. De foutgrens voorspelt een afname als $(\sin1)^n/(1 -
\sin1)$ --- ongeveer één cijfer per $13$ toetsaanslagen; de
a-posteriorigrens uit de weekendopgave van dit hoofdstuk (vraag
14) certificeert elke stap onderweg.
\end{example}

\section{Compactheid}

\begin{definition}\label{def:b2:metric:compact}
Een \omterm{def:b2:metric:def}{metrische ruimte} $X$ heet \emph{compact}\index{compacte
metrische ruimte} wanneer elke rij in $X$ een deelrij heeft die
\emph{in} $X$ convergeert (de eigenschap van
Bolzano--Weierstrass). Een deelverzameling heet compact wanneer zij
dat is met de geïnduceerde afstand.
\end{definition}

\begin{theorem}[Eerste eigenschappen]\label{thm:b2:metric:compactprops}
\begin{enumerate}
  \item Een \omterm{def:b2:metric:compact}{compacte} deelverzameling is gesloten en begrensd; een
        gesloten deelverzameling van een \omterm{def:b2:metric:compact}{compacte} ruimte is
        \omterm{def:b2:metric:compact}{compact}.
  \item In $\R^n$ geldt ook het omgekeerde: \omterm{def:b2:metric:compact}{compact} $\iff$ gesloten
        en begrensd.
  \item Een \omterm{def:b2:metric:continuity}{continu} beeld van een \omterm{def:b2:metric:compact}{compacte} ruimte is \omterm{def:b2:metric:compact}{compact}; een
        \omterm{def:b2:metric:continuity}{continue} reële functie op een niet-lege \omterm{def:b2:metric:compact}{compacte} ruimte is
        begrensd en neemt haar grenzen aan.
  \item (Heine) Een \omterm{def:b2:metric:continuity}{continue} afbeelding op een \omterm{def:b2:metric:compact}{compacte} ruimte is
        uniform \omterm{def:b2:metric:continuity}{continu}.
  \item Producten: zijn $X, Y$ \omterm{def:b2:metric:compact}{compact}, dan ook $X \times Y$ (met
        $d\bigl((x,y),(x',y')\bigr) = d(x,x') + d(y,y')$).
\end{enumerate}
\end{theorem}

\begin{proof}
(1) Dezelfde argumenten als op de rechte (volume van bachelorjaar
1): een rij die naar oneindig ontsnapt of buiten de verzameling
convergeert, heeft geen deelrij die binnenin convergeert; voor de
tweede bewering: extraheer in de omringende \omterm{def:b2:metric:compact}{compacte} ruimte en
gebruik de geslotenheid.

(2) Begrensde rijen in $\R^n$ hebben componentsgewijs convergente
deelrijen: extraheer op de eerste coördinaat
(Bolzano--Weierstrass op $\R$), daarna, uit die deelrij, op de
tweede, enzovoort ($n$ opeenvolgende extracties); de geslotenheid
houdt de limiet binnen.

(3) Gegeven $(f(x_n))$, extraheer $x_{\varphi(n)} \to x \in X$; de
\omterm{def:b2:metric:continuity}{continuïteit} geeft $f(x_{\varphi(n)}) \to f(x) \in f(X)$. Het reële
geval: de \omterm{def:b2:metric:compact}{compactheid} van $f(X) \subseteq \R$ maakt haar gesloten
en begrensd, en $\sup f(X) \in f(X)$ (het supremum van een
verzameling is haar aanhechtingspunt, en $f(X)$ is gesloten).

(4) Het bewijs uit bachelorjaar 1 gaat woordelijk over; hier is het
in metrische kleren. Stel dat $f \colon X \to Y$ \omterm{def:b2:metric:continuity}{continu} is op de
\omterm{def:b2:metric:compact}{compacte} $X$ maar niet uniform \omterm{def:b2:metric:continuity}{continu}: dan is er een $\varepsilon
> 0$ waarbij er voor elke $n$ punten zijn met
\[
d_X(x_n, y_n) \leq \frac{1}{n+1}
\qquad\text{en}\qquad
d_Y\bigl(f(x_n), f(y_n)\bigr) > \varepsilon .
\]
Extraheer $x_{\varphi(n)} \to a \in X$; dan geldt ook
$y_{\varphi(n)} \to a$ (de onderlinge afstanden gaan naar $0$). De
\omterm{def:b2:metric:continuity}{continuïteit} in $a$ stuurt beide beeldrijen naar $f(a)$, dus
$d_Y\bigl(f(x_{\varphi(n)}), f(y_{\varphi(n)})\bigr) \to 0$ --- in
strijd met de uniforme kloof $> \varepsilon$. De \omterm{def:b2:metric:compact}{compactheid}
leverde precies één ding: het verdichtingspunt $a$ waar de gewone
\omterm{def:b2:metric:continuity}{continuïteit} kan worden toegepast.

(5) Extraheer op de $X$-coördinaten en daarna nog eens op de
$Y$-coördinaten.
\end{proof}

\begin{example}[De stelling van Heine, met en zonder compactheid]\label{ex:b2:metric:heinework}
Op $\intcc{0}{1}$ is de functie $x \mapsto x^2$ uniform \omterm{def:b2:metric:continuity}{continu} ---
Heine zegt dat zonder enige berekening, maar de rechtstreekse
schatting is leerzaam:
\[
\abs{x^2 - y^2} = \abs{x + y}\,\abs{x - y} \leq 2\abs{x - y},
\]
dus voldoet $\delta = \varepsilon/2$ \emph{voor alle punten
tegelijk}. Op $\R$ is dezelfde functie niet uniform \omterm{def:b2:metric:continuity}{continu}: met
$x_n = n$ en $y_n = n + \frac1n$ gaat de kloof $\abs{x_n - y_n} =
\frac1n \to 0$ terwijl $\abs{x_n^2 - y_n^2} = 2 + \frac1{n^2} \geq
2$: geen enkele $\delta$ bedient $\varepsilon = 1$. Het mechanisme
is zichtbaar: de lokale Lipschitz-constante $\abs{x + y}$ is
begrensd op een \omterm{def:b2:metric:compact}{compacte} verzameling en onbegrensd op $\R$ --- de
stelling van Heine zegt precies dat \omterm{def:b2:metric:compact}{compactheid} zulke lokale
constanten uniform aftopt.
\end{example}

\begin{method}[Bewijzen dat een verzameling compact is]\label{met:b2:metric:compactness}
Drie routes, in volgorde van frequentie. (1) \emph{Herkenning in de
omringende ruimte:} ga in $\R^n$ (of in elke eindigdimensionale
genormeerde ruimte, \cref{ch:b2:nvs}) na dat de verzameling
gesloten is --- doorgaans als origineel,
\cref{ex:b2:metric:recognize} --- en begrensd. (2) \emph{Erven:}
een gesloten deelverzameling van een bekende \omterm{def:b2:metric:compact}{compacte} ruimte is
\omterm{def:b2:metric:compact}{compact}; een eindige vereniging of een product van \omterm{def:b2:metric:compact}{compacte}
verzamelingen is \omterm{def:b2:metric:compact}{compact}; een \omterm{def:b2:metric:continuity}{continu} beeld van een \omterm{def:b2:metric:compact}{compacte}
verzameling is \omterm{def:b2:metric:compact}{compact}. (3) \emph{Met blote handen:} extraheer uit
een willekeurige rij een convergente deelrij --- meestal door
opeenvolgende extracties, coördinaat voor coördinaat. Om
\emph{niet}-compactheid te bewijzen volstaat één getuige: een rij
zonder convergente deelrij, meestal punten op onderlinge afstand
$\geq \varepsilon$.
\end{method}

\begin{example}[Afstanden tussen verzamelingen: compactheid verdient haar loon]\label{ex:b2:metric:setdistance}
Zij $K$ \omterm{def:b2:metric:compact}{compact} en $F$ gesloten met $K \cap F = \emptyset$ in een
\omterm{def:b2:metric:def}{metrische ruimte}. Dan is
\[
d(K, F) = \inf\,\{d(x, y) : x \in K,\ y \in F\} > 0 :
\]
de functie $x \mapsto d(x, F)$ is \omterm{def:b2:metric:continuity}{continu}
(\cref{exo:b2:metric:11}) en positief op $K$ (uit $d(x, F) = 0$
zou $x \in \overline F = F$ volgen), dus neemt zij op de \omterm{def:b2:metric:compact}{compacte}
$K$ een positief minimum aan
(\cref{thm:b2:metric:compactprops}~(3)). De \omterm{def:b2:metric:compact}{compactheid} is geen
versiering: voor twee \emph{gesloten} verzamelingen kan het
infimum nul zijn zonder te worden aangenomen --- in $\R^2$ zijn de
hyperbool $F_1 = \{xy = 1\}$ en de as $F_2 = \{y = 0\}$ disjuncte
gesloten verzamelingen met $d(F_1, F_2) = 0$ (de punten $(n,
\frac1n)$ naderen de as). Ontsnappen naar oneindig is precies wat
\omterm{def:b2:metric:compact}{compactheid} verbiedt.
\end{example}

\begin{theorem}[Borel--Lebesgue]\label{thm:b2:metric:borellebesgue}
Een \omterm{def:b2:metric:def}{metrische ruimte} $X$ is \omterm{def:b2:metric:compact}{compact} dan en slechts dan als elke
overdekking van $X$ door \omterm{def:b2:metric:topology}{open} verzamelingen een \emph{eindige}
deeloverdekking heeft.\index{stelling van Borel--Lebesgue}
\end{theorem}

\begin{proof}
($\Leftarrow$) Stel dat $(x_n)$ geen convergente deelrij heeft. We
beweren dat elke $x \in X$ een bol $B(x, r_x)$ heeft die slechts
voor eindig veel indices $n$ het element $x_n$ bevat: anders zou
elke bol $B(x, \frac1{k+1})$ oneindig veel termen bevatten, en zou
het kiezen van indices
\[
\varphi(0) < \varphi(1) < \varphi(2) < \cdots
\quad\text{met}\quad
x_{\varphi(k)} \in B\Bigl(x, \frac{1}{k+1}\Bigr)
\]
(bij elke stap mogelijk juist omdat er oneindig veel kandidaten
overblijven) een deelrij bouwen die naar $x$ convergeert. De bollen
$B(x, r_x)$ overdekken $X$; zouden eindig veel van hen $X$
overdekken, dan zou de indexverzameling $\N$ een eindige vereniging
van eindige verzamelingen zijn: absurd.

($\Rightarrow$) Twee stappen. \emph{Lebesguegetal:} voor een \omterm{def:b2:metric:topology}{open}
overdekking $(U_i)$ van een \omterm{def:b2:metric:compact}{compacte} $X$ bestaat er een $\rho > 0$
zodanig dat elke bol met straal $\rho$ in een zekere $U_i$ ligt.
Anders kiezen we voor elke $n$ een $x_n$ waarvoor $B(x_n,
\frac{1}{n+1})$ in geen enkele $U_i$ ligt; extraheer
$x_{\varphi(n)} \to x \in U_{i_0} \supseteq B(x, r)$; voor grote
$n$ is $B(x_{\varphi(n)}, \frac{1}{\varphi(n)+1}) \subseteq B(x, r)
\subseteq U_{i_0}$: tegenspraak. \emph{Totale begrensdheid:} voor
elke $\rho > 0$ overdekken eindig veel bollen met straal $\rho$ de
ruimte $X$. Anders kiezen we inductief $x_{n+1}$ buiten $B(x_0,
\rho) \cup \dots \cup B(x_n, \rho)$: die rij heeft onderlinge
afstanden $\geq \rho$ en dus geen cauchydeelrij --- en dus geen
convergente deelrij: tegenspraak. Samengevoegd: overdek $X$ met
eindig veel bollen met straal $\rho$ (het lebesguegetal), elk
binnen een zekere $U_i$: dat is een eindige deeloverdekking.
\end{proof}

\begin{example}[Een $\varepsilon$-net, geteld]\label{ex:b2:metric:epsnet}
De totale begrensdheid (uit het bewijs van
\cref{thm:b2:metric:borellebesgue}) is op $\intcc{0}{1}$ heel
concreet: voor $\varepsilon > 0$ overdekken de $\lceil
\frac{1}{2\varepsilon}\rceil$ bollen met straal $\varepsilon$ rond
$\varepsilon, 3\varepsilon, 5\varepsilon, \dots$ het interval ---
ongeveer $\frac1{2\varepsilon}$ bollen, en geen enkele overdekking
kan het met minder dan $\frac{1}{2\varepsilon}$ ervan af (elke bol
dekt hoogstens lengte $2\varepsilon$). In $\intcc01^2$ wordt de
telling gekwadrateerd tot de orde $\varepsilon^{-2}$:
overdekkingsgetallen groeien als $\varepsilon^{-d}$ in dimensie
$d$ --- een kwantitatief gezicht van \omterm{def:b2:metric:compact}{compactheid}, en de reden dat
de oneindigdimensionale eenheidsbollen van \cref{ch:b2:nvs} (waar
helemaal geen eindig $\frac13$-net bestaat) niet \omterm{def:b2:metric:compact}{compact} kunnen
zijn.
\end{example}

\begin{example}[Compactheid aflezen aan overdekkingen]\label{ex:b2:metric:covers}
Het halfopen interval $\intoc{0}{1}$ wordt overdekt door de \omterm{def:b2:metric:topology}{open}
verzamelingen $U_n = \intoo{\frac1n}{2}$, $n \geq 1$; elke eindige
deelfamilie heeft een grootste index $N$ en mist
$\intoc{0}{\frac1N}$: geen eindige deeloverdekking, dus is
$\intoc{0}{1}$ niet \omterm{def:b2:metric:compact}{compact} --- wat de rijendefinitie ziet aan $x_n
= \frac1n$, waarvan de limiet $0$ ontsnapt. Anderzijds repareert
het toevoegen van het ene punt $0$ beide diagnoses tegelijk: op
$\intcc{0}{1}$ moet elke zulke overdekking een verzameling bevatten
die $0$ bevat, en die verzwelgt een heel beginstuk, waarna eindig
veel verzamelingen de rest afmaken. De twee talen van
\cref{thm:b2:metric:borellebesgue} falen of slagen altijd samen ---
overdekkingen zien ontsnapping precies waar rijen dat doen.
\end{example}

\begin{remark}[Vooruitblik binnen dit volume]\label{rem:b2:metric:perspectives}
Dit hoofdstuk is de dragende muur van het volume; let op waar elke
pijler gewicht draagt. \emph{\omterm{def:b2:metric:complete}{Volledigheid}}: het cauchycriterium
wordt de convergentietest voor reeksen in banachruimten
(\cref{ch:b2:series}), uniforme convergentie in
\cref{ch:b2:funcseq} is precies convergentie in het volledige
$\bigl(C, d_\infty\bigr)$, en Cauchy--Lipschitz
(\cref{ch:b2:diffeq}) is de vastepuntsstelling van Banach in het
kostuum van een integraalvergelijking. \emph{\omterm{def:b2:metric:compact}{Compactheid}}: zij
bewijst de equivalentie van normen (\cref{ch:b2:nvs}), het
aannemen van extrema bij de optimalisatie van
\cref{ch:b2:diffcalc}, en het bestaan van beste benaderingen (de
weekendopgave van \cref{ch:b2:nvs}). \emph{\omterm{def:b2:metric:connected}{Samenhang}}: zij
globaliseert lokale uitspraken --- de eenduidigheid van
oplossingen van differentiaalvergelijkingen, de
tussenwaardestelling op krommen (\cref{ch:b2:curves}), en de twee
componenten van $GL_n(\R)$ die de oriëntatietheorie
(\cref{ch:b2:multint}) uit elkaar zal houden.
\end{remark}

\begin{remark}[Klassieke valkuilen]\label{rem:b2:metric:pitfalls}
(i) ``Gesloten en begrensd impliceert \omterm{def:b2:metric:compact}{compact}'' is een stelling
over $\R^n$, niet over \omterm{def:b2:metric:def}{metrische ruimten}: een oneindige verzameling
met de discrete metriek is in zichzelf gesloten en begrensd en toch
niet \omterm{def:b2:metric:compact}{compact} (\cref{exo:b2:metric:4}), en de gesloten
eenheidsbol van $C(\intcc01)$ faalt eveneens (\cref{ch:b2:nvs}).
(ii) \omterm{def:b2:metric:complete}{Volledigheid} is een eigenschap van de \emph{afstand}, niet van
de topologie: $\R$ met $d(x,y) = \abs{\arctan x - \arctan y}$ heeft
dezelfde convergente rijen als gewoonlijk maar is onvolledig
(\cref{exo:b2:metric:1}). (iii) Een \omterm{def:b2:metric:continuity}{continue} bijectie hoeft geen
homeomorfisme te zijn --- de parametrisatie van de cirkel uit
\cref{exo:b2:metric:7}; \omterm{def:b2:metric:compact}{compactheid} van het vertrek repareert dat.
(iv) De stelling van Banach heeft $k < 1$ \emph{uniform} nodig: de
voorwaarde $d(f(x), f(y)) < d(x,y)$ alleen garandeert niets op een
\omterm{def:b2:metric:compact}{niet-compacte} ruimte (\cref{exo:b2:metric:5}). (v) \omterm{def:b2:metric:connected}{Samenhangend}
impliceert in het algemeen niet wegsamenhangend --- maar voor de
\omterm{def:b2:metric:topology}{open} deelverzamelingen van genormeerde ruimten die in dit boek
voorkomen, vallen beide samen (\cref{ch:b2:nvs}).
\end{remark}

\begin{remark}[Waar dit hoofdstuk wordt gebruikt]
Overal in de analysehelft. De \omterm{def:b2:metric:complete}{volledigheid} van $C(\intcc{a}{b})$
drijft de convergentiestellingen van \cref{ch:b2:funcseq} en de
theorie van Cauchy--Lipschitz in \cref{ch:b2:diffeq} aan (de
weekendopgave van dit hoofdstuk bewijst de lokale stelling van
Picard en Lindelöf al); \omterm{def:b2:metric:compact}{compactheid} geeft de equivalentie van
normen in eindige dimensie (\cref{ch:b2:nvs}) en het bestaan van
extrema in \cref{ch:b2:diffcalc}; \omterm{def:b2:metric:connected}{samenhang} ligt onder de
tussenwaardeargumenten van \cref{ch:b2:realfun} en onder het
globaal maken van de eenduidigheid voor
differentiaalvergelijkingen. In het volume van bachelorjaar 3
worden \omterm{def:b2:metric:compact}{compactheid} in functieruimten (de stelling van
Arzelà--Ascoli) en de categoriestelling van Baire
(\cref{exo:b2:metric:12} hier) dagelijks gereedschap.
\end{remark}

\begin{omfigure}
\begin{tikzpicture}[xscale=6.4, yscale=0.8]
  \draw[omDef, line width=3pt] (0,3) -- (1,3);
  \draw[omDef, line width=3pt] (0,2) -- (0.3333,2);
  \draw[omDef, line width=3pt] (0.6667,2) -- (1,2);
  \foreach \a in {0, 0.6667} {
    \draw[omDef, line width=3pt] (\a,1) -- (\a + 0.1111,1);
    \draw[omDef, line width=3pt]
      (\a + 0.2222,1) -- (\a + 0.3333,1);
  }
  \foreach \a in {0, 0.2222, 0.6667, 0.8889} {
    \draw[omThm, line width=3pt] (\a,0) -- (\a + 0.037,0);
    \draw[omThm, line width=3pt]
      (\a + 0.0741,0) -- (\a + 0.1111,0);
  }
  \node[left, font=\small] at (0,3) {$C_0$};
  \node[left, font=\small] at (0,2) {$C_1$};
  \node[left, font=\small] at (0,1) {$C_2$};
  \node[left, font=\small] at (0,0) {$C_3$};
\end{tikzpicture}

{\small De eerste stadia van de cantorverzameling
(\cref{exo:b2:metric:8}): elk niveau verwijdert het \omterm{def:b2:metric:topology}{open} middelste
derde van elk segment. De doorsnede $C = \bigcap_n C_n$ is \omterm{def:b2:metric:compact}{compact},
heeft leeg inwendige en lengte nul, en is toch \omterm{def:b2:structures:countable}{gelijkmachtig} met
$\R$ --- en zij keert terug als \emph{vast punt} van een contractie
op verzamelingen in de weekendopgave van dit hoofdstuk (vraag
22).}
\end{omfigure}

\section{Samenhang}

\begin{definition}\label{def:b2:metric:connected}
$X$ heet \emph{samenhangend}\index{samenhangende ruimte} wanneer
zij geen partitie in twee niet-lege \omterm{def:b2:metric:topology}{open} deelverzamelingen toelaat
--- gelijkwaardig: wanneer haar enige deelverzamelingen die zowel
\omterm{def:b2:metric:topology}{open} als gesloten zijn, $\emptyset$ en $X$ zijn. En $X$ heet
\emph{wegsamenhangend} wanneer elke twee punten door een \omterm{def:b2:metric:continuity}{continue}
afbeelding $\gamma \colon \intcc{0}{1} \to X$ worden verbonden.
\end{definition}

\begin{theorem}\label{thm:b2:metric:connectedness}
\begin{enumerate}
  \item De \omterm{def:b2:metric:connected}{samenhangende} deelverzamelingen van $\R$ zijn precies de
        intervallen.
  \item Een \omterm{def:b2:metric:continuity}{continu} beeld van een \omterm{def:b2:metric:connected}{samenhangende ruimte} is
        \omterm{def:b2:metric:connected}{samenhangend} --- waaruit de algemene tussenwaardestelling
        volgt: een \omterm{def:b2:metric:continuity}{continue} reële functie op een \omterm{def:b2:metric:connected}{samenhangende
        ruimte} neemt elke waarde aan tussen twee van haar waarden.
  \item Wegsamenhangend $\Rightarrow$ \omterm{def:b2:metric:connected}{samenhangend}. (Het omgekeerde
        faalt in het algemeen; het geldt voor \omterm{def:b2:metric:topology}{open}
        deelverzamelingen van genormeerde ruimten,
        \cref{ch:b2:nvs}.)
\end{enumerate}
\end{theorem}

\begin{proof}
(1) Een verzameling $A$ die geen interval is, mist een $z$ tussen
twee van haar punten: $A = (A \cap \intoo{-\infty}{z}) \cup (A \cap
\intoo{z}{+\infty})$ splitst haar in twee niet-lege (in $A$) \omterm{def:b2:metric:topology}{open}
stukken. Omgekeerd, zij $I$ een interval en $I = U \cup V$ een
partitie in niet-lege relatief \omterm{def:b2:metric:topology}{open} verzamelingen; kies $a \in U$
en $b \in V$, zeg $a < b$, en zet $s = \sup\,(U \cap
\intcc{a}{b})$, een punt van $\intcc{a}{b} \subseteq I$. Ligt $s
\in U$, dan is $s \neq b$, en legt de relatieve openheid van $U$
een heel interval rond $s$ (doorgesneden met $I$) binnen $U$ ---
zodat er punten van $U \cap \intcc{a}{b}$ groter dan $s$ zijn, in
strijd met het supremum. Ligt $s \in V$, dan legt de relatieve
openheid van $V$ een interval $\intoo{s - r}{s + r} \cap I$ binnen
$V$; maar het supremum is aanhechtingspunt van $U \cap
\intcc{a}{b}$, dat dat interval dus moet ontmoeten --- in strijd
met $U \cap V = \emptyset$. (Dit is het argument met open-en-gesloten
verzamelingen uit bachelorjaar 1 voor $\R$, uitgevoerd binnen $I$.)

(2) Splitst $f(X) = U' \cup V'$ in niet-lege relatief \omterm{def:b2:metric:topology}{open}
verzamelingen, dan splitst $X = f^{-1}(U') \cup f^{-1}(V')$ de
ruimte $X$ (\cref{thm:b2:metric:globalcontinuity}). De
tussenwaardestelling: $f(X) \subseteq \R$ is \omterm{def:b2:metric:connected}{samenhangend}, en dus
volgens (1) een interval.

(3) Stel $X = U \cup V$ met beide niet-leeg en \omterm{def:b2:metric:topology}{open}, en verbind $a
\in U$ met $b \in V$ door een weg $\gamma$: dan splitsen
$\gamma^{-1}(U)$ en $\gamma^{-1}(V)$ het interval $\intcc{0}{1}$,
in strijd met (1).
\end{proof}

\begin{example}\label{ex:b2:metric:glnr}
$GL_n(\R)$ is niet \omterm{def:b2:metric:connected}{samenhangend}: $\det$ is \omterm{def:b2:metric:continuity}{continu} (een veelterm in
de elementen) en surjectief op $\R^*$, dat niet \omterm{def:b2:metric:connected}{samenhangend} is; de
originelen van $\R_+^*$ en $\R_-^*$ splitsen $GL_n(\R)$. (Elk stuk
is in feite wegsamenhangend --- een aangename oefening, buiten wat
wij nodig hebben.) Daartegenover \emph{is} $GL_n(\C)$
wegsamenhangend: \cref{exo:b2:metric:10}.
\end{example}

\begin{example}[Een vast punt uit samenhang alleen]\label{ex:b2:metric:ivtfixed}
Elke \omterm{def:b2:metric:continuity}{continue} $f \colon \intcc01 \to \intcc01$ heeft een vast punt
--- zonder contractiehypothese, zonder iteratie. Beschouw $g(x) =
f(x) - x$, \omterm{def:b2:metric:continuity}{continu} op het \omterm{def:b2:metric:connected}{samenhangende} $\intcc01$:
\[
g(0) = f(0) \geq 0,
\qquad
g(1) = f(1) - 1 \leq 0 ,
\]
en de tussenwaardestelling
(\cref{thm:b2:metric:connectedness}~(2)) levert een nulpunt van
$g$, dat wil zeggen een vast punt van $f$. Vergelijk dit met
Banach (\cref{thm:b2:metric:banach}): hier is het bestaan
topologisch en gratis, maar de eenduidigheid en het algoritme zijn
verloren --- bij $f = \mathrm{id}$ ligt elk punt vast, en de
iteratie van een niet-samentrekkende $f$ kan eeuwig blijven
rondcirkelen. De twee vastepuntsstellingen van dit hoofdstuk
beantwoorden verschillende vragen in verschillende munteenheden.
\end{example}

\begin{example}[$\R$ en $\R^2$ zijn niet homeomorf]\label{ex:b2:metric:rnotr2}
\omterm{def:b2:metric:connected}{Samenhang} is een topologische vingerafdruk. Stel dat $h \colon \R^2
\to \R$ een homeomorfisme was (een \omterm{def:b2:metric:continuity}{continue} bijectie met \omterm{def:b2:metric:continuity}{continue}
inverse). Verwijder één punt $a \in \R^2$: de beperking $h \colon
\R^2\setminus\{a\} \to \R\setminus\{h(a)\}$ is nog steeds een
homeomorfisme. Maar $\R^2$ min één punt is wegsamenhangend ---
verbind twee punten door een segment, en maak een omweg langs een
tweede segment via een hulppunt als $a$ het rechtstreekse segment
blokkeert --- en dus \omterm{def:b2:metric:connected}{samenhangend}
(\cref{thm:b2:metric:connectedness}~(3)); terwijl $\R$ min één
punt in twee niet-lege \omterm{def:b2:metric:topology}{open} halfrechten uiteenvalt: niet
\omterm{def:b2:metric:connected}{samenhangend}. \omterm{def:b2:metric:connected}{Samenhang} blijft onder \omterm{def:b2:metric:continuity}{continue} afbeeldingen behouden:
tegenspraak. Het vlak en de rechte zijn werkelijk verschillend
\emph{als topologische ruimten} --- een feit dat kardinaliteit
alleen (bijecties in de trant van \cref{exo:b2:structures:3}
bestaan wel degelijk!) te grof is om te zien.
\end{example}

\section{Oefeningen}

\begin{exercise}[$\star$]\label{exo:b2:metric:1}
Ga op $\R$ na dat $\delta(x, y) = \min(1, \abs{x - y})$ en $d(x,y)
= \abs{\arctan x - \arctan y}$ afstanden zijn. Welke rijen
convergeren voor elk van beide? Is $(\R, d)$ \omterm{def:b2:metric:complete}{volledig}?
\end{exercise}

\begin{exercise}[$\star$]\label{exo:b2:metric:2}
Bewijs in een \omterm{def:b2:metric:def}{metrische ruimte} dat een convergente rij een
cauchyrij en begrensd is, en dat een cauchyrij met een convergente
deelrij convergeert. Leid daaruit opnieuw af dat \omterm{def:b2:metric:compact}{compacte metrische
ruimten} \omterm{def:b2:metric:complete}{volledig} zijn.
\end{exercise}

\begin{exercise}[$\star$]\label{exo:b2:metric:3}
Bereken in $\bigl(C(\intcc{0}{1}), d_\infty\bigr)$ de afstand
tussen $f(x) = x$ en $g(x) = x^2$; beschrijf de gesloten bol
$\overline B(0, 1)$; en bewijs dat de verzameling $\{f : f(0) =
0\}$ gesloten is terwijl $\{f : f(0) > 0\}$ \omterm{def:b2:metric:topology}{open} is.
\end{exercise}

\begin{exercise}[$\star\star$]\label{exo:b2:metric:4}
Bewijs dat de discrete \omterm{def:b2:metric:def}{metrische ruimte} $X$ (een willekeurige
verzameling) \omterm{def:b2:metric:complete}{volledig} is, en dat zij \omterm{def:b2:metric:compact}{compact} is dan en slechts dan
als $X$ eindig is. Welke deelverzamelingen zijn \omterm{def:b2:metric:connected}{samenhangend}?
\end{exercise}

\begin{exercise}[$\star\star$]\label{exo:b2:metric:5}
Zij $X$ \omterm{def:b2:metric:compact}{compact} en $f \colon X \to X$ met
\[
d\bigl(f(x), f(y)\bigr) < d(x, y) \quad \text{voor alle } x \neq y
.
\]
Bewijs dat $f$ precies één vast punt heeft \emph{(minimaliseer $x
\mapsto d(x, f(x))$)}, en geef een voorbeeld op $X =
\intco{1}{+\infty}$ (niet \omterm{def:b2:metric:compact}{compact}) zonder vast punt.
\end{exercise}

\begin{exercise}[$\star\star$]\label{exo:b2:metric:6}
(Geneste \omterm{def:b2:metric:compact}{compacte} verzamelingen) Zij $(K_n)$ een dalende rij van
niet-lege \omterm{def:b2:metric:compact}{compacte} deelverzamelingen van een \omterm{def:b2:metric:def}{metrische ruimte}.
Bewijs dat $\bigcap_n K_n \neq \emptyset$ \emph{(kies $x_n \in K_n$
en extraheer)}. Toon met een voorbeeld aan dat niet-lege geneste
\emph{gesloten} verzamelingen in $\R$ een lege doorsnede kunnen
hebben.
\end{exercise}

\begin{exercise}[$\star\star$]\label{exo:b2:metric:7}
Zij $K$ \omterm{def:b2:metric:compact}{compact} en $f \colon K \to Y$ \omterm{def:b2:metric:continuity}{continu} en bijectief. Bewijs
dat $f^{-1}$ \omterm{def:b2:metric:continuity}{continu} is \emph{(gebruik gesloten verzamelingen:
\cref{thm:b2:metric:globalcontinuity} en
\cref{thm:b2:metric:compactprops})}. Geef een tegenvoorbeeld zonder
\omterm{def:b2:metric:compact}{compactheid} ($\gamma(t) = (\cos t, \sin t)$ op
$\intco{0}{2\pi}$).
\end{exercise}

\begin{exercise}[$\star\star$]\label{exo:b2:metric:8}
De \emph{cantorverzameling} $C$ wordt uit $\intcc{0}{1}$ verkregen
door herhaaldelijk de \omterm{def:b2:metric:topology}{open} middelste derden te verwijderen. Bewijs
dat $C$ \omterm{def:b2:metric:compact}{compact} is, een leeg inwendige heeft, en oneindig is ---
sterker nog: \omterm{def:b2:structures:countable}{gelijkmachtig} met $\{0,1\}^{\N}$ \emph{(ternaire
ontwikkelingen met cijfers $0,2$; \cref{exo:b2:structures:3})}.
\end{exercise}

\begin{exercise}[$\star\star\star$]\label{exo:b2:metric:9}
Zij $X$ een \emph{\omterm{def:b2:metric:compact}{compacte}} \omterm{def:b2:metric:def}{metrische ruimte} en $f \colon X \to X$
een isometrie: $d(f(x), f(y)) = d(x, y)$. Bewijs dat $f$ surjectief
is. \emph{Aanwijzing: is $a \notin f(X)$, dan is $\varepsilon =
d(a, f(X)) > 0$ (waarom?); bestudeer de baan $a, f(a), f^2(a),
\dots$ en toon aan dat haar punten paarsgewijs minstens
$\varepsilon$ uit elkaar liggen --- in strijd met de \omterm{def:b2:metric:compact}{compactheid}.}
\end{exercise}

\begin{exercise}[$\star\star\star$]\label{exo:b2:metric:10}
Bewijs dat $GL_n(\C)$ wegsamenhangend is. \emph{Aanwijzing: zijn
$A, B$ inverteerbaar, beschouw dan $p(z) = \det\bigl((1 - z)A +
zB\bigr)$ voor $z \in \C$: een veelterm in $z$ die niet identiek
nul is, en die dus eindig veel nulpunten heeft; verbind $0$ met $1$
in $\C$ door een weg die ze vermijdt.}
\end{exercise}

\begin{exercise}[$\star\star$]\label{exo:b2:metric:11}
Zet voor $\emptyset \neq A \subseteq X$ de waarde $d(x, A) =
\inf_{a \in A} d(x, a)$. Bewijs dat $x \mapsto d(x, A)$
$1$-Lipschitz is, dat $d(x, A) = 0$ dan en slechts dan als $x \in
\overline A$, en dat voor disjuncte niet-lege \emph{gesloten}
verzamelingen $A, B$ de functie
\[
\varphi(x) = \frac{d(x, A)}{d(x, A) + d(x, B)}
\]
welgedefinieerd en \omterm{def:b2:metric:continuity}{continu} is, precies op $A$ gelijk aan $0$ en
precies op $B$ gelijk aan $1$ --- een \omterm{def:b2:metric:continuity}{continue} ``schakelaar'' die
elke twee disjuncte gesloten verzamelingen scheidt.
\end{exercise}

\begin{exercise}[$\star\star\star$]\label{exo:b2:metric:12}
(Baire) Zij $X$ een \omterm{def:b2:metric:complete}{volledige metrische ruimte} en $(U_n)_{n\geq1}$
een rij dichte \omterm{def:b2:metric:topology}{open} deelverzamelingen. Bewijs dat $\bigcap_n U_n$
dicht ligt in $X$ \emph{(bouw binnen een willekeurige bol geneste
gesloten bollen $\overline B(x_n, r_n) \subseteq U_n$ met $r_n \to
0$ en gebruik de \omterm{def:b2:metric:complete}{volledigheid})}. Leid af dat $\R$ geen \omterm{def:b2:structures:countable}{aftelbare}
vereniging is van gesloten verzamelingen met leeg inwendige, en
vind --- opnieuw --- dat $\R$ overaftelbaar is.
\end{exercise}

\section{Probleem: Picard-iteratie}

\omterm{def:b2:metric:complete}{Volledigheid} plus contractie is een oplosmachine: voer haar een
vergelijking die als vastepuntsprobleem is geschreven, en zij levert
bestaan, eenduidigheid, een algoritme en foutbalken. Deze
weekendopgave laat de machine op volle kracht draaien op de
vergelijking $y' = f(t, y)$: we bewijzen de lokale \emph{stelling
van Picard en Lindelöf} (het niet-lineaire hart van de theorie van
Cauchy--Lipschitz in \cref{ch:b2:diffeq}), zien elke hypothese haar
loon verdienen aan de hand van tegenvoorbeelden, en oogsten
zuiver metrische opbrengsten --- de \omterm{def:b2:metric:continuity}{continue} afhankelijkheid van de
gegevens, de vergelijking van Kepler, en de zelfgelijkvormigheid
van de cantorverzameling.

\begin{omfigure}
\begin{tikzpicture}
\begin{axis}[omaxis, width=8.6cm, height=5.4cm,
    xmin=-0.05, xmax=1.45, ymin=0, ymax=4.2,
    xtick={0,1}, ytick={1,2,3}]
  \addplot[omDef!45, thick, domain=0:1.4, samples=2] {1};
  \addplot[omDef!65, thick, domain=0:1.4, samples=2] {1 + x};
  \addplot[omDef, thick, domain=0:1.4, samples=60]
    {1 + x + x^2/2};
  \addplot[omThm, very thick, domain=0:1.4, samples=120]
    {exp(x)};
  \node[font=\small, omDef] at (axis cs:1.28,0.8) {$y_0$};
  \node[font=\small, omDef] at (axis cs:1.28,2.15) {$y_1$};
  \node[font=\small, omDef] at (axis cs:1.28,2.95) {$y_2$};
  \node[font=\small, omThm] at (axis cs:1.17,3.85)
    {$\eu^{\,t}$};
\end{axis}
\end{tikzpicture}

{\small Picard-iteraties voor $y' = y$, $y(0) = 1$: elke gang door
$T(y)(t) = 1 + \int_0^t y$ voegt één taylorterm toe, en de
contractie perst de hele rij uniform op $\eu^{\,t}$.}
\end{omfigure}

\begin{problem}[{Weekendopgave --- de stelling van Picard en
Lindelöf}]
\label{pb:b2:metric:1}
Overal zijn $t_0 \in \R$, $y_0 \in \R$, $a, b > 0$, en is $f$ een
\omterm{def:b2:metric:continuity}{continue} functie op de rechthoek $R = \intcc{t_0 - a}{t_0 + a}
\times \intcc{y_0 - b}{y_0 + b}$, begrensd door $M = \sup_R\,\abs
f$ en \emph{$L$-Lipschitz in haar tweede veranderlijke}:
$\abs{f(t, y) - f(t, z)} \leq L\abs{y - z}$ zodra beide punten in
$R$ liggen. Zet
\[
h = \min\Bigl(a, \frac bM\Bigr) \quad (\text{met } h = a
\text{ als } M = 0), \qquad I = \intcc{t_0 - h}{t_0 + h}.
\]

\medskip
\noindent\textbf{Deel I --- Het volledige toneel.}
\begin{enumerate}
  \item Bewijs de twee uitspraken die in
        \cref{def:b2:metric:complete} zijn aangehaald: een
        gesloten deelverzameling van een \omterm{def:b2:metric:complete}{volledige metrische
        ruimte} is \omterm{def:b2:metric:complete}{volledig}, en een volledige deelverzameling van
        een willekeurige \omterm{def:b2:metric:def}{metrische ruimte} is gesloten. Leid af dat
        elke gesloten deelverzameling van $\bigl(C(I),
        d_\infty\bigr)$ een \omterm{def:b2:metric:complete}{volledige metrische ruimte} is.
  \item Toon aan dat de afbeelding $C(I) \to C(I)$, $y \mapsto
        \bigl(t \mapsto \int_{t_0}^{t}y(s)\,\dd s\bigr)$,
        $h$-Lipschitz is voor $d_\infty$.
  \item (Vaste punten bewegen minder dan de afbeeldingen) Zij $g
        \colon X \to X$ een $k$-contractie van een \omterm{def:b2:metric:def}{metrische
        ruimte} met vast punt $\ell_g$, en $\widetilde g \colon X
        \to X$ een \emph{willekeurige} afbeelding met een vast
        punt $\ell_{\widetilde g}$. Bewijs dat
        \[
        d(\ell_g, \ell_{\widetilde g}) \leq
        \frac{d\bigl(g(\ell_{\widetilde g}),
        \widetilde g(\ell_{\widetilde g})\bigr)}{1 - k}
        \leq \frac{\sup_{x \in X} d\bigl(g(x), \widetilde
        g(x)\bigr)}{1 - k}.
        \]
  \item (De iteratietruc) Zij $X$ \omterm{def:b2:metric:complete}{volledig} en niet leeg en $g
        \colon X \to X$ een afbeelding --- niet noodzakelijk
        \omterm{def:b2:metric:continuity}{continu} --- waarvan een zekere iteratie $g^m$ een
        $k$-contractie is. Bewijs dat $g$ precies één vast punt
        $\ell$ heeft en dat \emph{elke} baan $x_{n+1} = g(x_n)$
        naar $\ell$ convergeert. \emph{(Vaste punten van $g$ zijn
        vaste punten van $g^m$; omgekeerd is $g(\ell)$ een vast
        punt van $g^m$; splits de baan naar de restklassen modulo
        $m$.)}
\end{enumerate}

\medskip
\noindent\textbf{Deel II --- De stelling van Picard en Lindelöf.}
\begin{enumerate}[resume]
  \item Toon aan dat een functie $y \colon I \to \intcc{y_0 -
        b}{y_0 + b}$ van klasse $C^1$ is met $y(t_0) = y_0$ en $y'
        = f(t, y)$ op $I$ dan en slechts dan als zij \omterm{def:b2:metric:continuity}{continu} is en
        aan de integraalvergelijking
        \[
        y(t) = y_0 + \int_{t_0}^{t} f\bigl(s, y(s)\bigr)\dd s
        \qquad (t \in I)
        \]
        voldoet.
  \item Zij $X_h = \{y \in C(I) : \abs{y(t) - y_0} \leq b
        \text{ op } I\}$ en zij $T$ gedefinieerd door $T(y)(t) =
        y_0 + \int_{t_0}^{t}f(s, y(s))\dd s$. Toon aan dat $X_h$
        een niet-lege gesloten deelverzameling van $C(I)$ is en
        dus \omterm{def:b2:metric:complete}{volledig}, en dat $T$ de ruimte $X_h$ in $X_h$
        afbeeldt --- hier doet $h \leq b/M$ zijn werk.
  \item Toon aan dat $d_\infty\bigl(T(y), T(z)\bigr) \leq
        Lh\,d_\infty(y, z)$ op $X_h$: is $Lh < 1$, dan besluit de
        stelling van Banach al. Deze kleinheidsvoorwaarde
        verwijderen we hierna.
  \item Bewijs met inductie naar $n$ dat
        \[
        \abs{T^n(y)(t) - T^n(z)(t)} \leq
        \frac{\bigl(L\abs{t - t_0}\bigr)^n}{n!}\,
        d_\infty(y, z) \qquad (y, z \in X_h,\ t \in I),
        \]
        zodat een zekere iteratie van $T$ een contractie is.
        Besluit met vraag 4 (\emph{de stelling van Picard en
        Lindelöf}): het beginwaardeprobleem $y' = f(t,y)$,
        $y(t_0) = y_0$ heeft precies één oplossing op $I =
        \intcc{t_0 - h}{t_0 + h}$ met waarden in $\intcc{y_0 -
        b}{y_0 + b}$.
  \item Toon aan dat de beperking ``met waarden in $\intcc{y_0 -
        b}{y_0 + b}$'' automatisch is: elke oplossing van het
        beginwaardeprobleem die op $I$ is gedefinieerd en waarvan
        de grafiek in $R$ begint, blijft in $\intcc{y_0 - b}{y_0 +
        b}$ \emph{(beschouw het eerste uittredemoment en begrens
        $\abs{y(t) - y_0}$ door $M\abs{t - t_0}$)}. De
        eenduidigheid geldt dus onder alle oplossingen op $I$.
  \item Laat de machine draaien op $y' = y$, $y(0) = 1$, startend
        bij de constante $y^{(0)} \equiv 1$: bereken de
        Picard-iteraties $y^{(n)}$, herken ze, en beschrijf de
        convergentie.
\end{enumerate}

\medskip
\noindent\textbf{Deel III --- Elke hypothese verdient haar loon.}
\begin{enumerate}[resume]
  \item (Zonder \omterm{def:b2:metric:continuity}{Lipschitz} geen eenduidigheid) Ga voor $y' =
        2\sqrt{\abs y}$, $y(0) = 0$ na dat $y \equiv 0$ en, voor
        elke $c \geq 0$, de functie $y_c(t) = 0$ voor $t \leq c$
        en $y_c(t) = (t - c)^2$ voor $t > c$, alle
        $C^1$-oplossingen op $\R$ zijn. Waar precies is $y \mapsto
        2\sqrt{\abs y}$ niet \omterm{def:b2:metric:continuity}{Lipschitz}?
  \item (De lokaliteit is echt) Los $y' = y^2$, $y(0) = 1$
        expliciet op, geef het maximale bestaansinterval, en
        bereken de beste $h$ die de stelling over alle keuzen van
        de rechthoek kan certificeren ($a$ groot, $b$ vrij): toon
        aan dat $h_{\max} = \sup_{b>0} \frac{b}{(1+b)^2} =
        \frac14$, terwijl de werkelijke oplossing op
        $\intoo{-\infty}{1}$ leeft.
  \item (\omterm{def:b2:metric:complete}{Volledigheid} is geen decoratie) Zij op $X = \Q \cap
        \intcc{1}{2}$ met de gebruikelijke afstand $g(x) = \frac
        x2 + \frac1x$. Toon aan dat $g(X) \subseteq X$, dat $g$
        een $\frac12$-contractie is \emph{(middelwaarde\-ongelijkheid)},
        en dat $g$ geen vast punt in $X$
        heeft. Welke hypothese van de stelling van Banach faalt,
        en wat is het vaste punt in de vervollediging?
  \item (Foutbalken) Bewijs voor een $k$-contractie $g$ op een
        volledige ruimte de \emph{a-posteriori}schatting $d(x_n,
        \ell) \leq \frac{k}{1-k}\,d(x_n, x_{n-1})$. Hoeveel
        stappen eist de \emph{a-priori}grens $\frac{k^n}{1-k}d(x_1,
        x_0)$ voor de afbeelding van Heron $g(x) = \frac x2 +
        \frac 1x$ op $\intcc{1}{2}$ (vast punt $\sqrt2$), startend
        bij $x_0 = \frac32$, om een nauwkeurigheid $10^{-6}$ te
        halen, en hoeveel stappen volstaan in werkelijkheid?
        (Bereken $x_1, x_2, x_3$ en hun fouten; de
        contractiegrens is eerlijk maar pessimistisch --- Heron
        convergeert kwadratisch.)
\end{enumerate}

\medskip
\noindent\textbf{Deel IV --- \omterm{def:b2:metric:continuity}{Continue} afhankelijkheid.} In dit deel
is $Lh < 1$, zodat $T$ zelf een contractie op $X_h$ is (vraag 7);
schrijf $y[\,y_0\,]$ voor de oplossing met beginwaarde $y_0$.
\begin{enumerate}[resume]
  \item (Afhankelijkheid van de beginwaarde) Zij $z_0$ een andere
        beginwaarde met $\abs{z_0 - y_0}$ klein genoeg opdat beide
        problemen in de rechthoek passen. Bewijs met vraag 3 dat
        \[
        d_\infty\bigl(y[y_0], y[z_0]\bigr) \leq
        \frac{\abs{y_0 - z_0}}{1 - Lh}.
        \]
  \item (Lange intervallen door aaneenschakelen) Neem aan dat de
        oplossingen bestaan op een lang segment dat in $m$
        opeenvolgende stukken is verdeeld waarop telkens de
        vorige grens geldt met $Lh \leq \frac12$. Toon aan dat de
        afwijking per stuk met hoogstens een factor $2$ groeit,
        en dus $d_\infty \leq 2^m\abs{y_0 - z_0}$ in het geheel
        --- een grens die exponentieel in de lengte is, de
        discrete schaduw van de $\eu^{L\abs{t - t_0}}$ uit het
        lemma van Gronwall (\cref{ch:b2:diffeq}).
  \item (Afhankelijkheid van het veld) Zij $g$ een ander veld op
        $R$, eveneens $L$-Lipschitz in $y$, met $\sup_R \abs{f -
        g} \leq \varepsilon$. Bewijs dat de bijbehorende
        oplossingen voldoen aan $d_\infty \leq \frac{\varepsilon
        h}{1 - Lh}$: modelfouten planten zich lineair voort.
  \item (Parameters) Is een familie $f_\lambda$ van velden
        uniform $L$-Lipschitz in $y$ en geldt
        $\sup_R\abs{f_\lambda - f_\mu} \leq C\abs{\lambda -
        \mu}$, leid dan af dat $\lambda \mapsto y_\lambda$
        \omterm{def:b2:metric:continuity}{Lipschitz} is van de parameterruimte naar $\bigl(C(I),
        d_\infty\bigr)$.
  \item (Stelsels kosten niets) Leg uit waarom de Delen I, II en
        IV woordelijk gelden voor $y$ met waarden in $\R^n$
        (supafstanden gebouwd op elk van de afstanden uit
        \cref{ex:b2:metric:examples}), en bereken vervolgens alle
        Picard-iteraties voor het stelsel $y' = Ay$, $y(0) =
        (c_1, c_2)$ met $A = \begin{psmallmatrix} 0 & 1\\ 0 &
        0\end{psmallmatrix}$: toon aan dat de iteratie na één stap
        stationair wordt op de exacte oplossing.
\end{enumerate}

\medskip
\noindent\textbf{Deel V --- Metrische opbrengsten en synthese.}
\begin{enumerate}[resume]
  \item (Storing van de identiteit) Zij $X$ een \omterm{def:b2:metric:complete}{volledig} toneel
        van het type genormeerde ruimte: neem $X = C(I)$ of
        $\R^n$. Is $\eta \colon X \to X$ $k$-Lipschitz met $k <
        1$, bewijs dan dat $x \mapsto x + \eta(x)$ een bijectie
        van $X$ is waarvan de inverse
        $\frac1{1-k}$-Lipschitz is \emph{(pas voor elke $y$
        Banach toe op $x \mapsto y - \eta(x)$)}. Dit is het
        metrische hart van de stelling van de inverse functie
        (\cref{ch:b2:diffcalc}).
  \item (De vergelijking van Kepler) Bewijs voor $0 \leq e < 1$ en
        $m \in \R$ dat $x = m + e\sin x$ precies één oplossing
        heeft, dat de iteratie $x_{n+1} = m + e\sin x_n$ er
        vanuit elk beginpunt naartoe convergeert, en schat: hoeveel
        iteraties garanderen voor $e = \frac12$ en $m = 1$ een
        fout $\leq 10^{-3}$ volgens de a-priorigrens? (De
        oplossing is $x \approx 1.4987$.)
  \item (De cantorverzameling is een vast punt) Zij $S_1(x) =
        \frac x3$ en $S_2(x) = \frac x3 + \frac23$ op $\R$, en zij
        $C$ de cantorverzameling uit \cref{exo:b2:metric:8}.
        Bewijs dat $C = S_1(C) \cup S_2(C)$, en leg in één zin uit
        waarom geen \emph{andere} niet-lege \omterm{def:b2:metric:compact}{compacte} verzameling
        aan deze vergelijking voldoet (de afbeelding $A \mapsto
        S_1(A) \cup S_2(A)$ is een contractie voor een afstand
        tussen \omterm{def:b2:metric:compact}{compacte} verzamelingen --- de
        hausdorffafstand, in het volume van bachelorjaar 3
        eerlijk gemaakt).
  \item (\omterm{def:b2:metric:connected}{Samenhang} maakt de eenduidigheid globaal) Zij $f$ lokaal
        \omterm{def:b2:metric:continuity}{Lipschitz} in $y$ op een \omterm{def:b2:metric:topology}{open} verzameling, en zijn $y, z$
        twee oplossingen van $y' = f(t, y)$ op een gemeenschappelijk
        interval $J$ met $y(t_0) = z(t_0)$. Bewijs dat $y = z$ op
        $J$: toon aan dat $\{t \in J : y(t) = z(t)\}$ niet leeg
        is, gesloten in $J$ en \omterm{def:b2:metric:topology}{open} in $J$ (wegens de lokale
        eenduidigheid), en gebruik de \omterm{def:b2:metric:connected}{samenhang} van intervallen
        (\cref{thm:b2:metric:connectedness}).
  \item (Geen kleinheid voor lineaire vergelijkingen) Pas voor $y'
        = \alpha(t)y + \beta(t)$ met $\alpha, \beta$ \omterm{def:b2:metric:continuity}{continu} op
        een segment $\intcc{A}{B}$ vraag 8 zo aan dat de grens met
        de faculteit op het \emph{hele} segment geldt, zodat
        bestaan en eenduidigheid daar globaal zijn --- het
        scalaire geval van de stelling van Cauchy--Lipschitz uit
        \cref{ch:b2:diffeq}, zonder enige beperking op de lengte
        $B - A$.
  \item (Synthese) In telkens één zin: wat de \omterm{def:b2:metric:complete}{volledigheid}
        bijdroeg; wat de contractie bijdroeg; wat de iteratietruc
        opleverde ten opzichte van Banach zonder meer; waar de
        \omterm{def:b2:metric:connected}{samenhang} binnenkwam; en welk tegenvoorbeeld van Deel III
        welke hypothese bewaakt. Noem de topstelling, en zeg wat
        de plaats van de contractie inneemt zodra $f$ alleen maar
        \omterm{def:b2:metric:continuity}{continu} is (de stelling van Peano, via \omterm{def:b2:metric:compact}{compactheid} in
        functieruimten --- Arzelà--Ascoli in het volume van
        bachelorjaar 3).
\end{enumerate}
\end{problem}
